\documentclass[../../../include/open-logic-section]{subfiles}

\begin{document}

\olfileid[hi]{sth}{ord-arithmetic}{using-addition}
\olsection{क्रमसूचक योग का उपयोग}

क्रमसूचकों पर योग का उपयोग करके हम
\olref[spine][rank]{defnsetrank} के अर्थ में विभिन्न समुच्चयों की रैंक
स्पष्टतः परिकलित कर सकते हैं:

\begin{lem}\ollabel{rankcomputation}
यदि $\setrank{A} = \alpha$ और $\setrank{B} = \beta$, तो:
\begin{enumerate}
	\item\ollabel{exrankpow} $\setrank{\Pow{A}} = \alpha\ordplus 1$
	\item\ollabel{exrankpair} $\setrank{\{A, B\}} = \max(\alpha,
	\beta) \ordplus 1$
	\item\ollabel{exrankcup} $\setrank{A \cup B} = \max(\alpha,
	\beta)$
	\item\ollabel{exranktuple} $\setrank{\tuple{A,B}} = \max(\alpha,
	\beta) \ordplus  2$
	\item\ollabel{exranktimes} $\setrank{A \times B} \leq \max(\alpha,
	\beta) \ordplus  2$
	\item\ollabel{exrankunion} जब $\alpha$ रिक्त या सीमा हो, तब
	$\setrank{\bigcup A} = \alpha$; जब $\alpha = \gamma\ordplus 1$, तब
	$\setrank{\bigcup A} = \gamma$।
\end{enumerate}
\end{lem}

\begin{proof}
पूरे प्रमाण में हम \olref[spine][rank]{ranksupstrict} का बार-बार आह्वान करते
हैं।

\emph{\olref{exrankpow}.} यदि $x \subseteq A$, तो
$\setrank{x} \leq \setrank{A}$। अतः
$\setrank{\Pow{A}} \leq \alpha \ordplus 1$। विशेषतः चूँकि
$A \in \Pow{A}$, इसलिए $\setrank{\Pow{A}} = \alpha \ordplus 1$।

\emph{\olref{exrankpair}.} \olref[spine][rank]{ranksupstrict} से।

\emph{\olref{exrankcup}.} \olref[spine][rank]{ranksupstrict} से।

\emph{\olref{exranktuple}.} \olref{exrankpair} से, दो बार।

\emph{\olref{exranktimes}.} ध्यान दें कि
$A \times B \subseteq \Pow{\Pow{A \cup B}}$, और \olref{exranktuple} का
आह्वान करें।

\emph{\olref{exrankunion}.} यदि $\alpha = \gamma\ordplus 1$, तो कोई
$c \in A$ ऐसा है कि $\setrank{c} = \gamma$, और $A$ के किसी !!{element}
की रैंक इससे अधिक नहीं; अतः $\setrank{\bigcup A} = \gamma$। यदि $\alpha$
सीमा क्रमसूचक है, तो $A$ के ऐसे !!{element}s हैं जिनकी रैंक $\alpha$ के
मनमाने रूप से निकट (पर कड़ाई से छोटी) है; इसलिए $\bigcup A$ के भी ऐसे
!!{element}s हैं जिनकी रैंक $\alpha$ के मनमाने रूप से निकट (पर कड़ाई से
छोटी) है; अतः $\setrank{\bigcup A} = \alpha$।
\end{proof}
\noindent
\olref{exranktimes} में \emph{अ}समानता क्यों आती है, यह दिखाना अभ्यास के
लिए छोड़ते हैं।

\begin{prob}
ऐसे समुच्चय $A$ और $B$ दीजिए कि
$\setrank{A \times B}= \max(\setrank{A}, \setrank{B})$। ऐसे समुच्चय $A$
और $B$ दीजिए कि
$\setrank{A \times B}\max(\setrank{A}, \setrank{B}) \ordplus 2$। क्या कोई
अन्य रैंक संभव है?
\end{prob}

अब हम यह भी दिखा सकते हैं कि किसी क्रमसूचक को ``परिमित'' या ``अनंत'' कहने
के अर्थ की कई उचित धारणाएँ परस्पर संपाती हैं:

\begin{lem}\ollabel{ordinfinitycharacter}
किसी भी क्रमसूचक $\alpha$ के लिए निम्न समतुल्य हैं:
\begin{enumerate}
	\item\ollabel{ord:notinomega} $\alpha\notin \omega$, i.e.,
	$\alpha$ प्राकृतिक संख्या नहीं है;
	\item\ollabel{ord:omegaplus} $\omega \leq \alpha$ 
	\item\ollabel{ord:oneplus} $1 \ordplus  \alpha = \alpha$ 
	%\alpha \approx \alpha \ordplus 1$
	\item\ollabel{ord:plusone} $\alpha \approx \alpha\ordplus 1$,
	अर्थात $\alpha$ और $\alpha\ordplus 1$ समसंख्य हैं;
	\item\ollabel{ord:infinite} $\alpha$ डेडेकिंड-अनंत है।
	\end{enumerate}
\end{lem}
\noindent
अतः किसी क्रमसूचक के (अ)परिमित होने का अर्थ समझने के हमारे पास पाँच सिद्धतः
समतुल्य ढंग हैं।

\begin{proof}
\emph{\olref{ord:notinomega} $\Rightarrow$ \olref{ord:omegaplus}.}
त्रिकोटी से।

\emph{\olref{ord:omegaplus} $\Rightarrow$ \olref{ord:oneplus}.}
$\alpha \geq \omega$ को नियत करें। परिमितेतर आगमन से कोई लघुतम क्रमसूचक
$\gamma$ (संभवतः $0$) ऐसा है कि कोई सीमा क्रमसूचक $\beta$ है जिसके लिए
$\alpha = \beta \ordplus \gamma$। अब:
\[
	1 \ordplus \alpha =  
	1 \ordplus (\beta \ordplus \gamma) = 
	(1 \ordplus \beta) \ordplus \gamma =  
	\supstrict_{\delta < \beta} (1 \ordplus  \delta) \ordplus  \gamma = 
	\beta \ordplus  \gamma = 
	\alpha.
\]
\emph{\olref{ord:oneplus} $\Rightarrow$ \olref{ord:plusone}.} स्पष्टतः
कोई !!a{bijection}
$f \colon (\alpha \disjointsum 1) \to (1 \disjointsum \alpha)$ है। यदि
$1 \ordplus \alpha = \alpha$, तो कोई समाकृतिकता
$g \colon (1 \disjointsum \alpha) \to \alpha$ है। अब $\comp{f}{g}$ पर
विचार करें।

\emph{\olref{ord:plusone} $\Rightarrow$ \olref{ord:infinite}.} यदि
$\alpha \approx \alpha \ordplus 1$, तो कोई !!a{bijection}
$f \colon (\alpha \disjointsum 1) \to \alpha$ है। प्रत्येक
$\gamma < \alpha$ के लिए $g(\gamma) = f(\gamma, 0)$ परिभाषित करें; यह
!!{injection} $\alpha$ के डेडेकिंड-अनंत होने का साक्षी है, क्योंकि
$f(0,1) \in \alpha \setminus \ran{g}$।

\emph{\olref{ord:infinite} $\Rightarrow$ \olref{ord:notinomega}.} यह
\olref[z][infinity-again]{naturalnumbersarentinfinite} है।
\end{proof}

\end{document}
